\cotefile{157-4}
\scanpage{288}
\mspage{16}
% archive sheet number: 287

des types $\mathrm{Hot}_{\unsure{W|w}}$. Si on tient compte de la description de $I_W$ comme $I\times e_W$, la formule d'adjonction \unsure{inhabituelle} (2.8.3) équivaut, bien sûr, à la formule (2.8.4) -- elle s'en déduit en ``\unsure{élevant} à la puissance $I$''.

On peut renverser la vapeur, et dans une catégorie \unsure{générale} $\mathcal H$, \uline{définir} un foncteur $\pi_0:\mathcal H\longrightarrow\mathrm{Ens}$ par la formule (2.8.\deleted{\unsure{4}}5)
\[
\pi_0(X)\overset{\mathrm{def}}{=}\mathrm{Hom}(e_{\mathcal H},X)
\quad\inserted{(=\Gamma_{\mathcal H}(X),\ \text{en une notation familière})}
\]
où $e_{\mathcal H}$ est \deleted{\unsure{l'}}\inserted{un} objet final de $\mathcal H$. On aura donc, en posant $I_{\mathcal H}=I\times e_{\mathcal H}$, la formule
\[
\mathrm{Hom}(I_{\mathcal H},X)\overset{\sim}{\longrightarrow}\mathrm{Hom}(\deleted{\unsure{\pi_0}}I,\deleted{\unsure{I}}\pi_0(X)),
\]
(i.e. la formule (2.8.3)), disant que l'on a un couple $(I_{\mathcal H},\pi_0)$ de foncteurs adjoints.

Supposons d'autre part que pour tout ens. $I$, l'application canonique (d'adjonction)
\[
 I\longrightarrow\deleted{\unsure{\mathrm{Hom}}}\pi_0(I_{\mathcal H})=\mathrm{Hom}_{\mathcal H}(e_{\mathcal H},e_{\mathcal H}\times I)
\]
% margin note written sideways
\margin{\unsure{Ça a tendance à} être \unsure{archi-faux} si $\mathcal H$ est additive. Par contre, vrai si $\mathcal H$ est une catégorie à sommes distributives \unsure{p.r.|p/r} aux produits (voire universelles) dont l'objet final est \uline{connexe} \ldots\ $\rightarrow$}
soit \uline{bijective},\inserted{$\ast$ i.e. que $I\to I_W$ soit \uline{pl. fid.}} On définit alors, pour \unsure{tt} $I$, une application canonique
\diagram{\begin{tikzcd}[ampersand replacement=\&,column sep=small]
\mathrm{Hom}(X,I_{\mathcal H}) \arrow[rr] \arrow[dr,dashed,"{\pi_0^{X,I_{\mathcal H}}}"'] \& \& \mathrm{Hom}(\pi_0(X),I) \arrow[dl,dashed,"\sim"]\\
\& \mathrm{Hom}(\pi_0(X),\pi_0(I_{\mathcal H})) \& 
\end{tikzcd}.}

\deleted{\unsure{L}}\unsure{La} ``bi-adjonction'' des foncteurs $\pi_0$ et $\unsure{\mathcal H|I}\mapsto I_{\mathcal H}$ revient à dire que cette flèche est \deleted{\unsure{p}}\unsure{bijective},
\scanpage{289}
\mspage{15}
% archive sheet number: 288

surjectif, donc $\exists\,s$ tel que $u_2s$ soit homotope à $u_1$ dans $\mathrm{Cat}$, ce qui implique $u_2s=u_1$ dans $\mathrm{Hot}_{\unsure{W|w}}$, \unsure{qed|qed.}.

\uline{Commentaire final}

L'axiome de connexité\inserted{Loc 5} diffère de tous les autres axiomes Loc 1 à Loc 8 (\unsure{mais|même} à l'exception de Loc 8 bis d'accessibilité) par le fait que ce n'est pas un axiome de stabilité.

Les autres axiomes Loc 1 à Loc 8 (et aussi Loc 8 bis d'accessibilité) et -- les axiomes \unsure{ultérieurs} \unsure{à considérer|de continuité}$^{\ast}$, \unsure{sont} satisfaits si $W=\unsure{\mathrm{Fl}|\mathrm{El}}\,M$ -- l'axiome de connexité est donc le seul qui s'oppose à la \uline{trivialité} de la théorie homotopique définie par $W$. C'est dire qu'il ne peut être conséquence de l'\unsure{ens.|un} des autres.
% slanted
\margin{il n'y en a en fait qu'un, grâce à la ``\unsure{providentielle}'' \unsure{thèse|théorie} de K.\unsure{S.} Brown !}

Par contre, l'axiome $\mathrm{Loc}(4)$ des sommes directes n'est pas indépendant des autres. Signalons à ce propos (rappel)
\[
\left\{\begin{aligned}
&\mathrm{Loc}(4a_0)\ \text{est conséquence de l'axiome $\mathrm{Loc}(6a)$}\\
&\hspace{3em}\text{des carrés cocartésiens (forme \unsure{directe|ouverte})}\\
&\mathrm{Loc}(4\unsure{a})\ \text{est conséquence de $\mathrm{Loc}(6a)$ précédent,}\\
&\hspace{3em}\text{$+$ l'axiome des limites inductives}\\
&\hspace{3em}\text{filtrantes $\mathrm{Loc}(8)$}\\
&\mathrm{Loc}(4b)\ \text{est conséquence de l'axiome $\mathrm{Loc}(7\unsure{)}$}\\
&\hspace{3em}\text{des fibrations $+$ \deleted{\unsure{donc}}}\\
&\mathrm{Loc}(4)\ \text{est conséquence de $\mathrm{Loc}(6a,7,8)$.}\\
&\hspace{3em}\text{(un point pour une vache,}\\
&\hspace{6em}\text{il est vrai\ldots)}
\end{aligned}\right.
\]
\scanpage{290}
\mspage{16}
% archive sheet number: 289

\unsure{D.} \uline{Axiome} $\mathrm{Loc}(6)$ \uline{des carrés cocartésiens} (ou de ``Mayer-Vietoris''\deleted{\unsure{)}}) (ou des $W$-cofibrations)
% slanted
\margin{C'est le nom ``\uline{axiome des $W$-cofibrations}'' qui m'apparaît à présent comme le meilleur, \unsure{car c'est|en tout cas} \unsure{ce nom|celui} qui me paraît le mieux saisir \unsure{la signification|ce qui} \unsure{conceptuelle|est en jeu} \inserted{\unsure{de l'axiome|en profondeur}}}

J'en vois plusieurs variantes, et n'ai pas encore l'esprit totalement net sur \deleted{\unsure{leurs}} \deleted{\unsure{dépendan}}\unsure{implications} mutuelles. Je commence par la variante plus compliquée, mais dont on aura besoin en tous cas :

\uline{Loc 6} (Axiome \inserted{$W$-cofibrations, ou encore Loc 5, -- axiome} des homomorphismes de carrés cocartésiens). Considérons \inserted{dans $\mathrm{Cat}$} un diagramme commutatif cubique
\diagram{\begin{tikzcd}[ampersand replacement=\&,row sep=3pt,column sep=small]
X_0 \arrow[rr,"{\unsure{i_2|\varepsilon_2}}"] \arrow[dr,"{\unsure{i_1|\varepsilon_1}}"'] \arrow[dd,"f_0"'] \& \& X_2 \arrow[dr] \arrow[dd,"f_2",crossing over] \& \\
\& X_1 \arrow[rr] \arrow[dd,"f_1"'] \& \& X_3 \arrow[dd,"f_3"]\\
Y_0 \arrow[rr,crossing over] \arrow[dr] \& \& Y_2 \arrow[dr] \& \\
\& Y_1 \arrow[rr] \& \& Y_3 ,
\end{tikzcd}}
% slanted
\margin{Il est bon d'introduire des propriétés du diagramme \deleted{\unsure{Loc}} $1^\circ)\ f_0\in W$ \quad $2^\circ)\ f_1$ et $f_2\in W$ \quad $3^\circ)\ f_3\in W$ et de formuler les parties a) b) c) de Loc 6 ainsi \quad a) $1^\circ+2^\circ\Rightarrow3^\circ$ \quad b) $2^\circ+3^\circ\Rightarrow1^\circ$ \quad c) $1^\circ+3^\circ\Rightarrow2^\circ$}
homomorphisme $f=(f_0,f_1,f_2,f_3)$ d'un diagramme carré $X$ dans un autre $Y$.

On suppose que $X$ provient d'un objet $X_3$ de $\mathrm{Cat}$, \deleted{\unsure{muni de deux}}\unsure{de} deux ouverts $X_1$, $X_2$ de $X_3$, \inserted{\uline{recouvrant} $X_3$} et de leur intersection $X_0=X_1\cap X_2$, et de même pour $Y$. (Appelons \inserted{``\uline{couvrants}''} ces diagrammes dans $\mathrm{Cat}$, diagrammes $\mathrm{MV}$ -- ``\uline{Mayer-Vietoris}'')
% slanted
\margin{\uline{NB} On verra plus tard (3.3) que $\uline{\mathrm{Loc}}(6c)$ implique déjà $\mathrm{Loc}(6a,b)$ i.e. $\mathrm{Loc}(6)$ tout entier. Et si $W$ satisfait $\mathrm{Loc}(1\,\mathrm{bis})$ p.~1. \unsure{plus fort|plus fortiori}, alors $\unsure{\mathrm{Loc}(6_0c)}\unsure{\Longleftrightarrow|\Leftarrow|\Rightarrow}\mathrm{Loc}(6a,b)$.}

On \unsure{demande} ce qui suit :
\[
\begin{aligned}
\text{a)}\quad f_0,f_1,f_2\in W &\Longrightarrow f_3\in W\\
\text{b)}\quad f_1,f_2,f_3\in W &\Longrightarrow f_0\in W\\
\bigl[\text{c)}\quad f_0,f_3\in W &\Longrightarrow f_1,f_2\in W\bigr]
\end{aligned}
\]
\margin{NB Il semble que dans $\mathrm{Loc}\,6a$, il suffit de l'\unsure{énoncer|axiome} en supposant $f_0,f_2$ des \uline{iso}, cf p.~22}
\margin{\unsure{Je ne faudrait|Il ne faudrait} pas l'inclure dans $\mathrm{Loc}\,6$ \quad cf p.~25}
\scanpage{291}
\mspage{17}
% archive sheet number: 290

ou encore que la flèche $\pi_0^{X,I_{\mathcal H}}$ :
\[
\mathrm{Hom}(X,I_{\mathcal H})\longrightarrow\mathrm{Hom}(\Gamma_{\mathcal H}(X),\Gamma_{\mathcal H}(I_{\mathcal H})\simeq I)
\]
soit bijective.

\uline{Corollaire} 2.\unsure{10}. \inserted{\deleted{(\unsure{avec les notations (2.8.3)})}} Le foncteur $I\unsure{\longleftarrow|\longmapsto}I_{\unsure{W|w}}$ de $\mathrm{Ens}$ dans $\mathrm{Hot}_W$ est pleinement fidèle, i.e. l'hom d'adjonction
\[
\pi_0(I_W)\longrightarrow I
\]
est bijectif. (Ou encore, $\pi_0$ est un foncteur de localisation.)

En effet, cette bijectivité est \deleted{\unsure{totale|visible}} évidente quand on travaille dans $\mathrm{Cat}$, et s'en déduit aussitôt \deleted{\unsure{pour}} dans $\mathrm{Hot}_W$.

\uline{Démonstration de} 2.8. La formule d'adjonction (2.8.2) résulte de la formule correspondante pour $\mathrm{Cat}$, évidente, et du lemme \unsure{2.5|25}. sur les foncteurs adjoints entre catégories de modèles. L'adjonction dans l'autre sens se réduit à la formule (2.8.4).

On a deux flèches canoniques
\diagram{\begin{tikzcd}[ampersand replacement=\&]
\mathrm{Hom}_W(e_W,X)\arrow[r,shift left,"\alpha"] \& \pi_0(X)\arrow[l,shift left,"\beta"],
\end{tikzcd}}
où $\alpha$ est $\pi_0^{e_W,X}$, compte tenu que $\pi_0(e_W)=\pi_0(e_{\mathrm{Cat}})=e_{\mathrm{Ens}}$, et $\mathrm{Hom}(e_{\mathrm{Ens}},\pi_0(X))\simeq\pi_0(X)$.

La flèche $\beta$ est \deleted{\unsure{déduite de}} la flèche \unsure{compo|comp}
\[
\mathrm{Ob}\,X\xrightarrow{\sim}\mathrm{Hom}_{\mathrm{Cat}}(e_{\mathrm{Cat}},X)\xrightarrow{\mathrm{hot}_W}\mathrm{Hom}_W(e_W,X)
\]
\scanpage{292}
\mspage{18}
% archive sheet number: 291

en notant que cette flèche passe au quotient en $\pi_0(X)\unsure{-|\longrightarrow}\mathrm{Hom}_{\unsure{W|w}}(e_W,X)$, vu que deux applications homotopes de $e_{\mathrm{Cat}}$ dans $X$ sont égales \unsure{dans} $\mathrm{Hot}_W$. Je dis que \deleted{\unsure{ce}}
\[
\beta\alpha=\mathrm{id},\qquad\alpha\beta=\mathrm{id}.
\]

La relation $\alpha\beta=\mathrm{id}$ revient à dire que la \deleted{l'application \illegible{1}}\inserted{flèche} $i_x:e_{\mathrm{Cat}}\unsure{-|\longrightarrow}X$ définie par $x\in\mathrm{Ob}\,\deleted{\unsure{X}}$, \deleted{donne un} donne $\pi_0(i_x)\deleted{\unsure{\text{élément}}}\inserted{{:}e_{\mathrm{Ens}}\to\pi_0(X)}$ \deleted{égale à} dont la valeur est la classe de $x$ dans $\pi_0(X)$, c'est tautologique. Pour prouver que $\beta\alpha=\mathrm{id}$, il suffit \inserted{\unsure{alors|donc}} de prouver que $\beta$ est surjective, i.e. le

\uline{\unsure{Korollaire|Corollaire}} 2.12. \inserted{Pour tt $X$ dans $\mathrm{Cat}$,} L'application
\[
\underbrace{\mathrm{Hom}_{\mathrm{Cat}}(e_{\mathrm{Cat}},X)}_{\simeq\mathrm{Ob}\,X}
\longrightarrow\mathrm{Hom}_W(e_W,X)
\]
est \uline{surjective} \deleted{\illegible{1}} (moyennant l'axiome de connexité Loc 5).

\deleted{Cela résulte du}

\uline{Dém. de 2.12}. Si $X=\varnothing$, l'assertion revient à dire que $\mathrm{Hom}_W(e_W,\varnothing)=\varnothing$, \deleted{\unsure{Or si}} \deleted{Mais il} mais plus généralement on a le

\uline{Cor. 2.1\deleted{\unsure{2}}3.} (Moyennant Loc 5) Soit $X$ non vide dans $\mathrm{Cat}$, alors $\mathrm{Hom}_W(X,\varnothing)=\varnothing$. \inserted{(Donc $\varnothing$ est un objet initial \uline{strict} de $\mathrm{Hot}_W$.)}

En effet, si on avait $X\to\varnothing$ dans $\mathrm{Hot}_W$,
\scanpage{293}
\mspage{17}
% archive sheet number: 292

% slanted
\margin{Je vais désigner \deleted{\unsure{Loc|par}} par Loc 6 l'axiome Loc 6(a,b), sans inclure c.}

\uline{NB} \unsure{1} Je distingue ces conditions comme $\mathrm{Loc}(6a)$, \deleted{\unsure{et}} $\mathrm{Loc}(6b)$, \inserted{$\mathrm{Loc}(6c)$} Je présume qu'on n'aura jamais à faire usage \unsure{à|de} $\mathrm{Loc}(6b)$, sans supposer également $\mathrm{Loc}(6a)$, mais l'inverse n'est sûrement pas le cas. Ainsi $W_0=\{f\in\unsure{\mathrm{Fl}|\mathrm{Fl.}}\mathrm{Cat}\mid\pi_0(f)\ \text{bijectif}\}$ satisfait a), mais non b). Je présume (sans avoir vérifié) que c'est pareil pour les $W_n$, $n\in\unsure{N|\mathbb N}$, de la ``\uline{$n$-équivalence homotopique}''.

Il se pourrait que $\uline{Loc}(6a)$ et $\uline{Loc}(6a,b)$ soient \unsure{équivalentes|équivalents} à la forme plus faible, et plus simple, suivante. (Si oui, c'est sous la forme la plus simple, bien sûr, qu'il convient de les énoncer.)

\uline{Loc} \deleted{\unsure{6}}$6'$ (Axiome des carrés cocartésiens, \inserted{ou de Mayer-Vietoris}).

Soit
\diagram{\begin{tikzcd}[ampersand replacement=\&,column sep=small]
X_0 \arrow[rr,"i_2"] \arrow[dr,"i_1"'] \& \& X_2\arrow[dr,"j_1"] \&\\
\& X_1\arrow[rr,"j_2"'] \& \& X_3
\end{tikzcd}}
un diagramme \deleted{\unsure{de}} $\mathrm{MV}$ dans $\mathrm{Cat}$. Alors
\[
\begin{aligned}
\text{a)}\quad&\text{Si } i_1\in W,\quad j_1\in W.\\
\text{b)}\quad&\text{Si } j_1\in W,\quad i_1\in W.
\end{aligned}
\]
\scanpage{294}
\mspage{18}
% archive sheet number: 293

\uline{Proposition} 3.1.
\[
\begin{aligned}
\mathrm{Loc}(6,a)&\Longrightarrow\mathrm{Loc}(6',a)\\
\deleted{\unsure{\mathrm{Loc}(6,b)}}&\deleted{\Longrightarrow\mathrm{Loc}(6'b)\ \&\ \mathrm{Loc}(6'a)}\\
\mathrm{Loc}(6\unsure{,}b)&\Longrightarrow\mathrm{Loc}(6'b)\\
\mathrm{Loc}(6c)&\Longrightarrow\mathrm{Loc}(6'a,b)
\end{aligned}
\]
% slanted
\margin{NB on a même $\mathrm{Loc}\,6c\Longrightarrow\underbrace{\mathrm{Loc}\,6(\deleted{\unsure{c}}a,b)}_{\mathrm{Loc}\,6}$ \deleted{\unsure{$\Longrightarrow$}}}

Ça se voit sur les diagrammes cubiques \deleted{\unsure{S}} $Q$
\diagram{\begin{tikzcd}[ampersand replacement=\&,row sep=small,column sep=small]
X_0\arrow[rr,"i_2"]\arrow[dr,"i_1"'] \& \& X_2\arrow[dr,"j_1"] \&\\
\& X_1\deleted{\unsure{X_1}}\arrow[rr,"{\deleted{\unsure{f_2}}\,j_2}"] \& \& X_3\\
X_0\arrow[uu,"f_0=\mathrm{id}"]\arrow[rr,"i_2",crossing over]\arrow[dr,"\mathrm{id}"'] \& \& X_2\arrow[uu,"{\deleted{\unsure{\mathrm{id}}}f_2=\mathrm{id}}"',crossing over]\arrow[dr,no head,"\mathrm{id}"] \&\\
\& X_0\arrow[uu,"f_1=i_1"]\arrow[rr,"i_2"'] \& \& X_2\arrow[uu,"f_3=j_1"']
\end{tikzcd}\quad $(Q)$,}
\diagram{\begin{tikzcd}[ampersand replacement=\&,row sep=tiny,column sep=tiny]
X_0\arrow[rr,"i_2"]\arrow[dr,"i_1"']\arrow[dd,"i_1"'] \& \& X_2\arrow[dr,"j_1"]\arrow[dd,"j_1",crossing over] \&\\
\&X_1\arrow[rr,"j_2"]\arrow[dd,"\mathrm{id}"'] \& \& X_3\arrow[dd,"\mathrm{id}"]\\
X_1\arrow[rr,"j_2",crossing over]\arrow[dr,"\mathrm{id}"'] \& \& X_3\arrow[dr,"\mathrm{id}"] \&\\
\& X_1\arrow[rr,"j_2"'] \& \& X_3
\end{tikzcd}\quad $(Q')$}
où les $f_0,f_1,f_2,f_3$ de l'énoncé de $\mathrm{Loc}(6)$ deviennent $\mathrm{id}_{X_0},i_1,\mathrm{id}_{X_2},j_1$. $\mathrm{Loc}(6a)$ donne donc $i_1\in W\Longrightarrow j_1\in W$, \inserted{i.e. $\mathrm{Loc}(6'a)$} et $\mathrm{Loc}(6b)$ $j_1\in W\Longrightarrow i_1\in W$, \deleted{\unsure{d'où la}}\inserted{i.e. $\mathrm{Loc}(6'b)$}

\deleted{proposition 3.\unsure{1|2}.} Pour $\mathrm{Loc}(6b)\Longrightarrow\mathrm{Loc}(6',b)$, on utilise \inserted{(cf \uline{XII} p.~\unsure{78|77})} le cube $(Q')$
\diagram{\deleted{\begin{tikzcd}[ampersand replacement=\&,row sep=tiny,column sep=tiny]
\unsure{X_0}\arrow[rr,"{\unsure{i_2}}"]\arrow[dr,"{\unsure{i_1}}"]\arrow[dd,"{\unsure{i_2}}"] \& \& \unsure{X_2}\arrow[dr,"{\unsure{j_1}}"]\arrow[dd,"{\unsure{\mathrm{id}}}"] \&\\
\&\unsure{X_1}\arrow[rr,"{\unsure{j_2}}"]\arrow[dd,"{\unsure{j_2}}"] \& \& X_3\arrow[dd,"{\unsure{\mathrm{id}}}"]\\
\unsure{X_2}\arrow[rr,"{\unsure{\mathrm{id}}}"]\arrow[dr,"{\unsure{j_1}}"] \& \&\unsure{X_2}\arrow[dr,"{\unsure{j_1}}"] \&\\
\& \unsure{X_3}\arrow[rr,"{\unsure{\mathrm{id}}}"] \& \&\unsure{X_3}
\end{tikzcd}}}
, où $f_0=i_1$, $f_1=\mathrm{id}_{X_1}$, $f_{\unsure{2|1}}=j_1$, $f_3=\mathrm{id}_{X_3}$, \inserted{donc $j_1\in W\Longrightarrow i_1\in W$}
% slanted
\margin{\inserted{$\uline{Loc}(6b)\Rightarrow$} i.e. $\uline{Loc}(6'b)$, et si on a $\uline{Loc}(6c)$, on voit que $i_1\in W\Rightarrow j_1\in W$ i.e. $\mathrm{Loc}(6'a)$.}

\uline{Commentaire 3.\unsure{2|?}.} $\mathrm{Loc}(6a)$ m'apparaît comme le ``minimum vital'' pour établir l'existence d'une ``théorie'' des \uline{co}fibres $W$-homotopiques dans $\mathrm{Cat}$. Pour cela il faut en effet disposer d'un critère serviable assurant qu'une flèche $i:A\unsure{-|\longrightarrow}B$ dans $\mathrm{Cat}$ est dans $\mathrm{Cof}_{\unsure{W|w}}$.
\scanpage{295}
\mspage{19}
% archive sheet number: 294

On trouve, quand on dispose de $\mathrm{Loc}(6,a)$\inserted{(\unsure{sa forme|ses formes} \unsure{duale|duales}\inserted{i.e. pour les \uline{fermés}}, c'est kif-kif)}, que les \uline{immersions ouvertes de Dwyer} sont dans $\mathrm{Cof}_{\unsure{W|w}}$, \deleted{\unsure{d'où}} ce qui suffit pour avoir \uline{une \inserted{bonne} théorie} \uline{des cofibres} $W$-\uline{homotopiques}, \uline{dans \unsure{laquelle|lesquelles} les carrés} $\mathrm{MV}$ \uline{sont} $W$-\uline{cocartésiens}. (chose à laquelle on tient absolument).

Quant à $\mathrm{Loc}(6b)$, \deleted{\unsure{c'est}} si je dispose à la place de $\mathrm{Loc}(6'b)$, \inserted{\uline{et $\mathrm{Loc}(6a)$,}} je dois m'en tirer pour prouver $\mathrm{Loc}(6)$
\deleted{et \unsure{même}, pour \unsure{avoir} $\mathrm{Loc}(6)$ sous la forme en apparence plus forte de l'ancien axiome $W(7)$ (``axiome du cube'' ou ``lemme des cinq'', cf \uline{XII} p.~76, \unsure{i.e.|ici} aussi la forme}
\[
\deleted{\text{c)}\quad f_0,f_3\in W\Longrightarrow f_1,f_2\in W.}
\]

Dans XII Appendice (pages \unsure{265 ff|265 A}) je prouve, paraît-il, que (avec les notations actuelles) $\mathrm{Loc}(6')$ implique $\mathrm{Loc}(6)$, et $\unsure{=|\Longleftrightarrow}\mathrm{Loc}(6c)$ si $W$ stable par facteurs directs (condition $\mathrm{Loc}(1\,\mathrm{bis})$). Mais la démonstration est incomplète, puisque j'admets tacitement que les immersions de Dwyer sont des $\mathrm{Cof}_W$, chose que je n'avais établie que moyennant l'ax. $W(7)$, en fait moyennant $\mathrm{Loc}(6'a)$, ce qui est une partie de $W(7)$).
\scanpage{296}
\mspage{2\unsure{0|9}}
% archive sheet number: 295

Mais ici me vient l'idée du lemme de K.T. Brown, cité par \unsure{Baues}, que je vais redonner ici :

\uline{Lemme} 3.2.1. (K.T. Brown) Soit $(M,W)$ un localiseur ($W$ modérément saturé), muni d'un ens. $C\subset\mathrm{Fl}\,M$ satisfaisant les conditions :

1) $C$ stable par composition et par cochangement de base (NB on n'a \deleted{pas} à supposer l'existence des sommes amalgamées
\diagram{\begin{tikzcd}[ampersand replacement=\&,row sep=small,column sep=small]
X_0\arrow[rr,"i_1"]\arrow[dr,"i_2"']\& \&X_1\arrow[dr] \&\\
\& X_2\arrow[rr,no head] \& \&X_3
\end{tikzcd}}
que si $i_1\in C$), contient les iso.

2) $M$ contient un objet initial, $\varnothing_M$, et pour tout objet \inserted{$X$} de $M$, $\varnothing\longrightarrow X$ soit dans $C$. (NB 1) et 2) impliquent que $M$ est stable par sommes finies)

% A long arrow carries this item above the preceding condition.
\inserted{\deleted{\unsure{4}}3) Toute flèche $f$ de $M$ se factorise en $pi$, avec $i\in C$, $p\in W$.}

\deleted{3}4) $C\cap W$ est stable par cochangement de base, i.e. si
\diagram{\begin{tikzcd}[ampersand replacement=\&,row sep=small,column sep=small]
X_0\arrow[rr,no head,"i_1"]\arrow[dr,"{i_2\deleted{\unsure{i_2}}}"'] \& \&X_1\arrow[dr,"{\deleted{\unsure{j_1}}j_2}"] \&\\
\& X_2\arrow[rr,no head,"j_1"'] \& \&X_3
\end{tikzcd}}
est cocartésien et \deleted{\unsure{ons}} $i_1\in C\cap W$, alors $j_1\in C\cap W$. (NB $j_1\in C$ est déjà contenu dans 1), il reste \inserted{à exiger} $j_1\in W$)

Sous ces conditions, on a $\deleted{\unsure{W}}C\subset\mathrm{Cof}_{\unsure{W|w}}$.

On peut aussi énoncer le lemme ainsi

\uline{Cor.} 3.2.2. Supposons que $C$ satisfait 1, 2, 3. Alors \deleted{la condition 3) \uline{équivaut} à $C\subset\mathrm{Cof}_W$,}
\scanpage{297}
\mspage{2\deleted{\unsure{0}}1}
% archive sheet number: 296

\deleted{\unsure{On exprime|On suppose} \unsure{celle-ci|ceci} : \unsure{dans|étant donné} un carré cocartésien $(\ast)$, si $i_1\in C$, $i_2\in W$, alors $j_2\in W$.}

les conditions suivantes sur $C$ sont équivalentes

a) Condition \deleted{3}4) ci-dessus, i.e. stabilité de $C\cap W$ par cochangt de base

\[
\text{a$'$)}\quad C\cap W\subset W^{\unsure{\mathrm{univ}}}
\]

\[
\text{b)}\quad C\subset\mathrm{Cof}_{\unsure{W|w}}
\]

b$'$) Pour tt carré cocartésien
\diagram{\begin{tikzcd}[ampersand replacement=\&,column sep=small,row sep=small]
X_0\arrow[rr,"{\unsure{i|i_2}}"]\arrow[dr,"p"']\arrow[dr,phantom,shift left=4ex,"\mathrm{coc}"] \& \& X_1\arrow[dr,"q"] \&\\
\& X'_0\arrow[rr,no head,"j"] \& \& X'_1,
\end{tikzcd}}
avec $i\in C$, $p\in W$, on a $q\in W$.

\uline{Dém} \qquad On a
\diagram{\deleted{\begin{tikzcd}[ampersand replacement=\&]
a\arrow[r,Leftrightarrow] \& a'\\
b\arrow[r,Leftrightarrow]\arrow[u,Rightarrow] \& b'\arrow[u,Rightarrow]
\end{tikzcd}}}
\[
\begin{aligned}
b&\Longleftrightarrow b'\\
a&\Longleftrightarrow a'
\end{aligned}
\]
car $C$ stable par \unsure{coch.} de base. On a $b\Longrightarrow a'$ car $\mathrm{Cof}_W\cap W=W^{\mathrm{univ}}$. \inserted{grâce à l'hyp. 3$^\circ$)} On a $a\Longrightarrow b'$ par le lemme de Brown.

On va appliquer ceci au cas $M=\mathrm{Cat}$, $C=$ immersions ouvertes de Dwyer. Les conditions 1) 2) ne dépendent pas de $W$, elles sont vérifiées, et 3) aussi dès que $W$ est un localiseur fondamental. Quant à 4), c'est une conséquence de l'axiome $W(7\,\mathrm{bis})$, i.e. $\mathrm{Loc}(6')$ (on a
\scanpage{298}
\mspage{22}
% archive sheet number: 297

besoin apparemment de $\mathrm{Loc}(6'a\ \text{\uuline{et}}\ b)$.) On en conclut $C\subset\mathrm{Cof}_{\unsure{W|w}}$ -- c'est merveilleux ! (Mais je ne sais pas prouver le lemme \ldots)

Cela légitime la démonstration faite en \uline{XII} loc. cit., et on trouve la

\uline{Proposition} 3.3. La condition $\mathrm{Loc}(6')$ (axiome des carrés cocartésiens, forme forte\inserted{a \& b}) \deleted{implique} \uline{équivaut} à $\mathrm{Loc}(6.a,b)$ (axiome des hom. de carrés cocartésiens)\inserted{et est \uline{impliquée} par $\mathrm{Loc}(6c)$.} \deleted{\unsure{forme} \illegible{2}}
\deleted{(qui implique (\unsure{déjà, th?} 18) \unsure{d'après} prop. 3.1.)}
\inserted{Si $W$ satisfait à $\mathrm{Loc}(1\,\mathrm{bis})$, i.e. stable par fact. dir. alors $\mathrm{Loc}(6')$ équivaut et à $\mathrm{Loc}(6abc)$ \deleted{et à $\mathrm{Loc}(6c)$} \inserted{$\mathrm{Loc}(6,c)$.}}
% slanted
\margin{\uline{\unsure{Commentaire|Corollaire}}}

3.\deleted{\unsure{3}}4. Ainsi, on trouve une théorie des cofibres $W$-homotopiques \uline{soit} sous l'hypothèse $\mathrm{Loc}(6')$, (équivalente à $\mathrm{Loc}(6)$\inserted{plus \unsure{compliquée|compliqué}}), soit sous l'hypothèse \deleted{plus} faible\inserted{et $+$ compliquée} $\mathrm{Loc}(6a)$.

Dans le premier cas, c'est une théorie \inserted{de nature} très particulière, puisque les carrés $W$-cocartésiens sont ``réversibles'', (i.e. on a \deleted{\unsure{implicat}} non seulement $i_1\in W\Longrightarrow j_1\in W$, mais aussi l'implication inverse).

Il reste donc la question si $\mathrm{Loc}(6'a)$ \deleted{est équivalente} suffit pour une théorie des cofibres \inserted{$W$-}homotopiques, ou ce qui revient sûrement au $=$, si ça implique déjà $\mathrm{Loc}(6a)$ (donc lui est équivalent), C'est \unsure{entièrement|ultérieurement} suspendu à la \uline{question} :
\scanpage{299}
\mspage{23}
% archive sheet number: 298

\margin{?}

l'axiome $\mathrm{Loc}(6'a)$ (axiome des carrés cocartésiens, forme directe) implique-t-il que \inserted{l'ens.} des immersions ouv. de Dwyer qui sont \unsure{des|dans} $W$, \deleted{\unsure{est}} (imm. de Dwyer $W$-triviales) est stable par cochangement de base ?
En fait, on voit qu'il suffirait pour cela de vérifier l'axiome du cube (ou des hom de carrés cocartésiens) \inserted{forme directe $\mathrm{Loc}(6a)$}, dans le cas où $f_0$ et $f_2$ sont des identités, donc à prouver qu'alors $f_1\in W\Longrightarrow f_3\in W$. (ce qui donnerait directement $\mathrm{Dw}\subset\mathrm{Cof}_{\unsure{W|w}}$). Mais je ne vois aucun moyen de le prouver\ldots

Supposons maintenant que la condition $\mathrm{Loc}(6a)$ est vérifiée, donc qu'on dispose d'une ``bonne théorie des cofibres $W$-homotopiques'', donc aussi d'une théorie des carrés $W$-cocartésiens\inserted{(cf \uline{XII} p.~258--262)}. Il y a deux questions qui se posent, concernant la validité, avec ces hypothèses plus faibles, \unsure{des développements|du développement} de \uline{XII} :

$1^\circ$) Soit un carré commutatif
\diagram{\begin{tikzcd}[ampersand replacement=\&,column sep=small,row sep=small]
X_0\arrow[rr]\arrow[dr] \& \& X_2\arrow[dr] \&\\
\& X_1\arrow[rr]\&\& X_3
\end{tikzcd}}
de $\mathrm{Cat}$, est-il vrai que ce carré est $W$-cocartésien ssi la flèche canonique
\diagram{\begin{tikzcd}[ampersand replacement=\&,column sep=tiny,row sep=tiny]
\&\& X_1\&\&\\
\displaystyle\int \& X_0\arrow[ur]\arrow[dr] \& \& {}\arrow[r,no head] \& X_3\\
\&\&X_2\&\&
\end{tikzcd}}
est dans $W$ ?
% slanted
\margin{Mieux, \unsure{en tous les cas|ici dans les cas}, quand on a un diagramme
\diagram{\begin{tikzcd}[ampersand replacement=\&,row sep=0pt,column sep=tiny]
\&X_1\\X_0\arrow[ur]\arrow[dr]\&\\\&X_2
\end{tikzcd}}
, son intégrale
\diagram{\begin{tikzcd}[ampersand replacement=\&,row sep=0pt,column sep=tiny]
\&\&{}\\\displaystyle\int\&{}\arrow[ur]\arrow[dr]\&\\\&\&{}
\end{tikzcd}}
est une ``somme amalg. \unsure{entre|ent.} $W$-homotopique'' \unsure{dudit} diagramme (moyennant $\mathrm{Loc}(6a)$ \unsure{sans plus}, et $=$ seulement $\mathrm{Loc}(6'a)$)}
\scanpage{300}
\mspage{2\deleted{\unsure{3}}4}
% archive sheet number: 299

$2^\circ$) Est-il encore vrai que si $A$ est dans $\mathrm{Cat}$, $u:F\hookrightarrow G$ un \deleted{flèche}\inserted{monomorphisme} de $\widehat A$, alors $u\in\mathrm{Cof}_{\unsure{W|w}_A}$, et si $u\in W$, $u\in W_A^{\unsure{\mathrm{univ}}}$, donc a-t-on
\[
\mathrm{Mono}_{\widehat A}\subset\mathrm{Cof}_{W_A},\quad
\mathrm{Mono}_{\widehat A}\cap W_A\subset(W_A)^{\mathrm{univ}}\ ?
\]
% slanted
\margin{NB $\mathrm{Loc}(6'a)$ suffit (et $\mathrm{Loc}(6a)$)}

Mais \unsure{$1^\circ$)|ceci} est prouvé \unsure{ds|dans} \uline{XII} prop 20 p 262 -- je viens de vérifier que la démonstr. est valable \deleted{dans} les hypothèses actuelles.$^{\ast}$ \unsure{pourvu|pour} que je vérifie que le cor. 3 p.~149 \uline{XII} est valable, i.e.
% slanted
\margin{Mais elle utilise de plus $\mathrm{Loc}(3\,\mathrm{bis})$. Mais avec \unsure{cette|la} démonstration \uline{plus simple} que je viens de \unsure{rajouter|regarder} (p.~264), on peut s'en passer (et on obtient une vision b. plus complète de la situation)}

$1^\circ$ bis) Si
\diagram{\begin{tikzcd}[ampersand replacement=\&,row sep=small,column sep=small]
X_0\arrow[r,no head]\arrow[d] \& X_1\arrow[d]\\
X_2\arrow[r,no head]\&X
\end{tikzcd}}
est un carré $\mathrm{MV}$ dans $\mathrm{Cat}$, alors
\diagram{\begin{tikzcd}[ampersand replacement=\&,row sep=tiny,column sep=tiny]
\&\&X_1\&\&\\
\displaystyle\int\&X_0\arrow[ur]\arrow[dr]\&\&{}\arrow[r]\&X\\
\&\&X_2\&\&
\end{tikzcd}}
est dans $W$. (donc, moyennant $1^\circ$), le carré est $W$-cocartésien).

Pour le vérifier, je suis amené à vérifier que le th. 4 de loc. cit. (p.~\unsure{146|145}) est encore valable. / je viens de le faire, et ai annoté en conséquence le th. 4 et ses corollaires (certains résultats \inserted{(cor. 4. et cor 5)} utilisent $W(7\,\unsure{bis})$ dans \deleted{toute} sa force, i.e. $\mathrm{Loc}(6')$ i.e. $\mathrm{Loc}(6,a\ \text{\uline{et}}\ b)$. Mais \deleted{\unsure{ces résultats}} cor. 3 du th. 4 est OK.

Je viens aussi de vérifier la validité de $2^\circ$), en vérifiant la dém. du cor. du th 5 (\uline{XII} p.~168) et en l'annotant en conséquence.
\scanpage{301}
\mspage{2\deleted{\unsure{4}}5}
% archive sheet number: 300

\deleted{3.5. Je résume les principaux résultats obtenus}
\deleted{\inserted{$\mathrm{Loc}\,6\unsure{c}\Longleftrightarrow$}}
\[\tag*{\deleted{(3.5.1)}}\deleted{\mathrm{Loc}\,6c\Longrightarrow\mathrm{Loc}\,6(a,\deleted{b})\deleted{\unsure{ab}}\Longleftrightarrow\mathrm{Loc}\,6'}\]
\[\tag*{\deleted{(3.5.2)}}\begin{aligned}&\deleted{\text{Si } W\text{ satisfait } \mathrm{Loc}\,1\,\mathrm{bis}\ \deleted{\unsure{\text{strict}}}\ (W=W^{\unsure{\natural|\sharp}}),\ \text{on a}}\\&\deleted{\mathrm{Loc}\,6\unsure{a|c}\Longleftrightarrow}\end{aligned}\]

3.5. Résumons : on a
\[\tag*{(3.5.1)}\mathrm{Loc}\,6c\ \deleted{\unsure{\Longleftrightarrow\mathrm{Loc}\,6}}\Longrightarrow\mathrm{Loc}\,6\deleted{\unsure{(a,b)}}\Longleftrightarrow\mathrm{Loc}\,6'\]
et
\[\tag*{(3.5.2)}\text{si }W=W^{\natural},\ \text{on a}\quad\mathrm{Loc}\,6c\Longleftrightarrow\mathrm{Loc}\,6\Longleftrightarrow\mathrm{Loc}\,6'\deleted{\unsure{(a,b)}}\]
\deleted{$\unsure{\mathrm{Loc}\,6'}$.}

\deleted{J'ignore si $\mathrm{Loc}\,6'a$ implique $\mathrm{Loc}\,6a$}

D'autre part
\[\tag*{(3.5.3)}\mathrm{Loc}\,6a\Longrightarrow\mathrm{Loc}\,6'a,\]
et j'ignore si \deleted{$6a$ implique} l'implication inverse
\[\mathrm{Loc}\,6'a\overset{?}{\Longrightarrow}\mathrm{Loc}\,6a\]
% slanted
\margin{C'est le cas en présence de $\mathrm{Loc}(8,9)$}
est valable. Pour établir $\mathrm{Loc}\,6a$, il suffit de le faire en \unsure{tous cas|le cas} quand $f_0$ et $f_2$ sont des iso.

L'axiome $\mathrm{Loc}\,6a$ semble le ``minimum vital'' pour que $W$ donne lieu à une ``bonne théorie des cofibres $W$-homotopiques'', ou encore des carrés $W$-cocartésiens. De façon plus précise, j'ai envie de prouver ceci :

\uline{Proposition} 3.6. L'axiome $\mathrm{Loc}\,6a$ équivaut à la conjonction des deux propriétés suivantes

$\mathrm{Loc}\,6\,\mathrm{A}$ : \deleted{(\unsure{appelons})} $\mathrm{Cat}$ a ``assez de $W$-cofibrations''
\scanpage{302}
\mspage{26}
% archive sheet number: 301

i.e. \unsure{toute|tout} \unsure{flèche|flèches} $f$ de $\mathrm{Cat}$ se factorise en $f=pi$, où $i\in\mathrm{Cof}_{\unsure{W|w}}$ et $p\in W$. (Cela signifie aussi que $\mathrm{Cat}$ munie de $(W,\mathrm{Cof}_W)$ est une catégorie à cofibrations de K. S. Brown--Anderson, et bien des résultats de la théorie des \unsure{Brown-Quillen|Brown-oreilles} sont applicables\deleted{\unsure{.)}}. En particulier, on a une bonne théorie des carrés $W$-cocartésiens.)

\uline{$\mathrm{Loc}\,6\,\mathrm B$} Si \deleted{$\unsure{X}$} dans $\mathrm{Cat}$ est réunion de deux ouverts $U,V$, d'intersection $w$, alors le carré
% slanted
\margin{En fait, $\mathrm{Loc}\,6\,\mathrm B$ est conséquence de $\mathrm{Loc}\,6\,\mathrm A$, \deleted{\unsure{(qui} je prouve} \inserted{laquelle lui} est \uline{équivalente} à $\mathrm{Loc}\,6a$ ? Cf esquisse de démonstration \deleted{page} des 4.4 bis pages 66--68.}
\diagram{\begin{tikzcd}[ampersand replacement=\&]
w\arrow[r]\arrow[d]\&U\arrow[d]\\
V\arrow[r,dash]\&X
\end{tikzcd}}
est $W$-cocartésien.

D'autre part, la condition $\mathrm{Loc}\,6\,\mathrm A$ \deleted{\unsure{est cette}} \inserted{peut être} \deleted{équivalente à} \inserted{remplacée par} la condition (en apparence plus faible)

\uline{$\mathrm{Loc}\,6\,\mathrm A'$} Tout morphisme de Dwyer est dans $\mathrm{Cof}_W$.

\uline{Corollaire} 3.7. La condition $\mathrm{Loc}\,6\deleted{\unsure{(ab)}}$ \inserted{i.e $\mathrm{Loc}\,6(a,b)$} (ou encore la condition équivalente $\mathrm{Loc}\,6'$) équivaut à la conjonction des conditions $\mathrm{Loc}\,6\,\mathrm A$ \inserted{(ou $\mathrm{Loc}\,6\,\mathrm A'$)} et $\mathrm{Loc}\,6\,\mathrm B$ ci-dessus, et de la condition

\uline{$\mathrm{Loc}\,6\,\mathrm C$} : Pour \deleted{tout} carré \inserted{$W$-}cocartésien
\diagram{\begin{tikzcd}[ampersand replacement=\&]
X_0\arrow[r,"\unsure{i|j}"]\arrow[d]\&X_1\arrow[d]\\
X_{\unsure{2|1}}\arrow[r,"j"]\&\unsure{X|X^{\prime}}
\end{tikzcd}}
dans $\mathrm{Cat}$, $j\in W\Longrightarrow i\in W$.
